\subsection{Differentiation and Basic Theorems}

% **Role:** one-variable differentiation as local linearization and a source of global estimates.
%
% Derivatives, mean-value theorems, inverse differentiation, and Taylor approximation.

This subsection develops the derivative as a relative-domain limit and first-order approximation.\\
It then derives algebraic, mean-value, inverse, higher-order, and vector-valued consequences.

\begin{definition} \label{def:derivative-real-function} \hlt{Derivative of a Real Function}\\
Let $X\subseteq\R$, let $f:X\to\R$, and let $x_0\in X^{\operatorname{acc}}$. Define $\phi_{x_0}:X\setminus\{x_0\}\to\R$ by $\phi_{x_0}(x)\defeq\frac{f(x)-f(x_0)}{x-x_0}$.\\
If limit exists, define $f'(x_0)\defeq\lim\limits_{x\to x_0}\phi_{x_0}(x)$ as \hlt{derivative at $x_0$ relative to $X$}. Map $f$ is then \hlt{differentiable at $x_0$ relative to $X$}.\\
For $E\subseteq X^{\operatorname{acc}}$, $f$ is \hlt{differentiable on $E$ relative to $X$} $\Longleftrightarrow(\forall x\in E)\,f'(x)$ exists relative to $X$.
\end{definition}

\begin{perspective} \label{per:one-variable-domain-of-differentiability} \hlt{Derivative Depends on the Domain}\\
The punctured approach set $X\setminus\{x_0\}$ determines the derivative in Definition \ref{def:derivative-real-function}.\\
Interior interval points produce two-sided derivatives, whereas interval endpoints produce one-sided derivatives.\\
The limit-point hypothesis excludes isolated points.
\end{perspective}

\noindent Figure \ref{fig:one-variable-domain-restriction} contrasts an isolated point with a half-line endpoint.\\
Only the endpoint admits a punctured approach within its displayed domain.

\begin{illustrationfigure}[H]
  \centering
  \begin{tikzpicture}[
  scale=0.92,
  >=Stealth,
  every node/.style={font=\scriptsize},
  axis/.style={->, gray!70, line width=0.45pt},
  domainline/.style={blue!65!black, line width=1.0pt},
  approach/.style={->, orange!80!black, line width=0.7pt},
  point/.style={circle, fill=black, inner sep=1.15pt},
  guide/.style={gray!55, dashed, line width=0.45pt},
  labelbox/.style={fill=white, fill opacity=0.9, text opacity=1, inner sep=1.2pt}
]
  \begin{scope}
    \draw[axis] (0,0) -- (4.4,0) node[right] {$\R$};
    \foreach \x in {0.75,2.2,3.6} {
      \node[point] at (\x,0) {};
    }
    \draw[guide] (1.62,-0.26) rectangle (2.78,0.26);
    \node[above=2pt] at (2.2,0) {$x_0$};
  \end{scope}

  \begin{scope}[xshift=5.3cm]
    \draw[axis] (0,0) -- (4.4,0) node[right] {$\R$};
    \draw[domainline] (0.75,0) -- (4.05,0);
    \node[point] at (0.75,0) {};
    \node[above=2pt] at (0.75,0) {$x_0$};
    \node[labelbox, above=11pt] at (2.45,0.26) {$X=[x_0,\infty)$};
    \draw[approach] (3.55,0.32) -- (2.75,0.07);
    \draw[approach] (2.7,0.32) -- (1.95,0.07);
    \draw[approach] (1.9,0.32) -- (1.05,0.07);
  \end{scope}
\end{tikzpicture}

  \caption{Left: $x_0$ is isolated. Right: $X=[x_0,\infty)$ permits right-sided approach to $x_0$.}
  \label{fig:one-variable-domain-restriction}
\end{illustrationfigure}

\noindent Figure \ref{fig:derivative-secant-to-tangent} represents a difference quotient by a secant slope. Its limiting value is the tangent slope $f'(x)$.

\begin{illustrationfigure}[H]
  \centering
  \begin{tikzpicture}[
  scale=0.37,
  >=Stealth,
  every node/.style={font=\scriptsize},
  axis/.style={->, gray!70, line width=0.45pt},
  curve/.style={blue!70!black, line width=0.8pt},
  secant/.style={orange!80!black, line width=0.7pt},
  tangent/.style={green!45!black, line width=0.7pt},
  guide/.style={gray!45, dashed, line width=0.4pt},
  point/.style={circle, fill=gray!85, inner sep=1.4pt},
  labelbox/.style={fill=white, fill opacity=0.88, text opacity=1, inner sep=1.2pt}
]
  \draw[axis] (-0.625,0) -- (12.4,0) node[right] {$t$};
  \draw[axis] (0,-0.5) -- (0,8.7) node[above] {$y$};

  \draw[curve, domain=0.25:9.5, samples=90]
    plot (\x,{0.875*(\x/2.5-1.2)*(\x/2.5-1.2)+1.75});

  \coordinate (X) at (3.5,1.785);
  \coordinate (T) at (8.0,5.25);

  \draw[secant] (1.875,0.525) -- (9.3,6.25);
  \draw[tangent] (1.125,1.4525) -- (9.875,2.6775);

  \draw[guide] (3.5,0) -- (X);
  \draw[guide] (8.0,0) -- (T);
  \node[below] at (3.5,0) {$x$};
  \node[below] at (8.0,0) {$t$};

  \node[point] at (X) {};
  \node[point] at (T) {};
  \node[labelbox, anchor=south east] at (3.5,2.45)
    {\textcolor{blue!70!black}{$(x,f(x))$}};
  \node[labelbox, anchor=west] at (8.55,5.05)
    {\textcolor{blue!70!black}{$(t,f(t))$}};

  \node[labelbox, align=center] at (5.75,5.25)
    {\textcolor{orange!75!black}{$\frac{f(t)-f(x)}{t-x}$}};
  \node[labelbox] at (8.65,3.2)
    {\textcolor{green!45!black}{$f'(x)$}};
\end{tikzpicture}

  \caption{The secant through $(x,f(x))$ and $(t,f(t))$ has slope $\frac{f(t)-f(x)}{t-x}$; its limit is $f'(x)$.}
  \label{fig:derivative-secant-to-tangent}
\end{illustrationfigure}

\begin{proposition} \label{prop:one-variable-first-order-approximation} \hlt{First-Order Approximation Form}\\
Let $X\subseteq\R$, let $x_0\in X^{\operatorname{acc}}$, let $f:X\to\R$, and let $L\in\R$. Set $R_L(x)\defeq f(x)-f(x_0)-L(x-x_0)$.\\
Then $f'(x_0)=L$ $\Longleftrightarrow(\forall\varepsilon>0)(\exists\delta>0)(\forall x\in X\setminus\{x_0\})\,[\abs{x-x_0}<\delta\Longrightarrow\abs{R_L(x)}<\varepsilon\abs{x-x_0}]$.
\begin{proof}
For $x\in X\setminus\{x_0\}$, $\frac{\abs{R_L(x)}}{\abs{x-x_0}}=\abs{\frac{f(x)-f(x_0)}{x-x_0}-L}$.
\begin{enumerate}[label=\roman*.]
\setlength{\itemsep}{0pt}
\item[$(\Rightarrow)$] Assume $f'(x_0)=L$, $\varepsilon>0$. On some punctured $\delta$-neighborhood, the right-hand side is less than $\varepsilon$ by Definition \ref{def:derivative-real-function}.\\
Multiplication by $\abs{x-x_0}$ proves the estimate.
\item[$(\Leftarrow)$] Assume the estimate, and $\varepsilon>0$. Choose $\delta>0$ satisfying the estimate by the hypothesis.\\
Division by $\abs{x-x_0}$ makes right-hand side less than $\varepsilon$ on same punctured neighborhood. Hence $f'(x_0)=L$ by Definition \ref{def:derivative-real-function}.
\end{enumerate}
\hlt{Structural Commentary}: The two directions rescale the same remainder by $\abs{x-x_0}$.\\
\hlt{Key Insight}: The normalized first-order error is the derivative-limit error.
\end{proof}
\end{proposition}

\begin{theorem} \label{thm:differentiable-implies-continuous} \hlt{Differentiability Implies Continuity}\\
Let $X\subseteq\R$, let $x_0\in X^{\operatorname{acc}}$, and let $f:X\to\R$. If $f$ is differentiable at $x_0$ relative to $X$, then $f$ is continuous at $x_0$.
\begin{proof}
For $t\in X\setminus\{x_0\}$, $f(t)-f(x_0)=\frac{f(t)-f(x_0)}{t-x_0}(t-x_0)$. Factors tend to $f'(x_0)$ and $0$, so $f(t)-f(x_0)\to0$ by Theorem \ref{thm:function-limit-algebra}(ii).\\
Thus $f(t)\to f(x_0)$ by Theorem \ref{thm:function-limit-algebra}(i), and continuity holds by Theorem \ref{thm:continuity-as-limit-at-limit-point}.\\
\hlt{Structural Commentary}: Limit algebra transfers difference-quotient convergence to the function increment.\\
\hlt{Key Insight}: A finite quotient multiplied by a vanishing increment tends to zero.
\end{proof}
\end{theorem}

\begin{theorem} \label{thm:derivative-algebra-rules} \hlt{Algebra Rules for Derivatives}\\
Let $X\subseteq\R$, let $x_0\in X^{\operatorname{acc}}$, and let $f,g:X\to\R$ be differentiable at $x_0$ relative to $X$.
\begin{enumerate}[label=(\roman*)]
\setlength{\itemsep}{0pt}
\item $(f+g)'(x_0)=f'(x_0)+g'(x_0)$.
\item $(fg)'(x_0)=f'(x_0)g(x_0)+f(x_0)g'(x_0)$.
\item If $g(x_0)\neq0$, then $\left(\frac{f}{g}\right)'(x_0)=\frac{g(x_0)f'(x_0)-f(x_0)g'(x_0)}{g(x_0)^2}$.
\end{enumerate}
\begin{proof}
For $h\in\{f,g\}$ and $t\in X\setminus\{x_0\}$, set $Q_h(t)\defeq\frac{h(t)-h(x_0)}{t-x_0}$; then $Q_h(t)\to h'(x_0)$ as $t\to x_0$.
\begin{enumerate}[label=(\roman*)]
\setlength{\itemsep}{0pt}
\item Since $Q_{f+g}=Q_f+Q_g$, the sum rule holds by Theorem \ref{thm:function-limit-algebra}(i).
\item The identity $Q_{fg}(t)=f(t)Q_g(t)+g(x_0)Q_f(t)$ holds.\\
The limit $f(t)\to f(x_0)$ holds by Theorem \ref{thm:differentiable-implies-continuous}, so the product rule follows by Theorem \ref{thm:function-limit-algebra}(i),(ii).
\item Assume $g(x_0)\neq0$. Then $g(t)\neq0$ near $x_0$ by Theorem \ref{thm:differentiable-implies-continuous} and the assumption.\\
On this punctured neighborhood, $Q_{f/g}(t)=\frac{g(x_0)Q_f(t)-f(x_0)Q_g(t)}{g(t)g(x_0)}$. The quotient rule follows by Theorem \ref{thm:function-limit-algebra}(i)--(iii).
\end{enumerate}
\hlt{Structural Commentary}: Difference-quotient identities reduce all three claims to limit algebra and continuity.\\
\hlt{Key Insight}: Expand increments before passing to their limits.
\end{proof}
\end{theorem}

\begin{theorem} \label{thm:chain-rule-one-variable} \hlt{Chain Rule}\\
Let $X,Y\subseteq\R$, $f:X\to Y$, $g:Y\to\R$, $x_0\in X^{\operatorname{acc}}$. Set $y_0\defeq f(x_0)$ and assume $y_0\in Y^{\operatorname{acc}}$.\\
If $f'(x_0)$ and $g'(y_0)$ exist relative to $X$ and $Y$, then $(g\circ f)'(x_0)=g'(y_0)f'(x_0)$ relative to $X$.
\begin{proof}
Define $u:X\to\R$ by $u(x_0)\defeq0$ and $u(t)\defeq\frac{f(t)-f(x_0)}{t-x_0}-f'(x_0)$ for $t\neq x_0$.\\
Define $v:Y\to\R$ by $v(y_0)\defeq0$ and $v(s)\defeq\frac{g(s)-g(y_0)}{s-y_0}-g'(y_0)$ for $s\neq y_0$. Then $u(t)\to0$ and $v(s)\to0$ by Definition \ref{def:derivative-real-function}.\\
The limit $f(t)\to y_0$ holds by Theorem \ref{thm:differentiable-implies-continuous}. The definition $v(y_0)=0$ and punctured limit $v(s)\to0$ therefore imply $v(f(t))\to0$.\\
For $t\in X\setminus\{x_0\}$, $\frac{g(f(t))-g(f(x_0))}{t-x_0}=\bigl(g'(y_0)+v(f(t))\bigr)\bigl(f'(x_0)+u(t)\bigr)$. Right-hand side tends to $g'(y_0)f'(x_0)$ by Theorem \ref{thm:function-limit-algebra}(i),(ii).\\
\hlt{Structural Commentary}: Composed first-order expansions reduce the derivative to a product of two limit factors.\\
\hlt{Key Insight}: The linear terms multiply while both normalized remainders vanish.
\end{proof}
\end{theorem}

\begin{example} \label{ex:oscillatory-derivative-examples} \hlt{Oscillation at a Point}\\
Assume the usual derivative of $\sin$, and define $f_1,f_2:\R\to\R$ by
\begin{align*}
f_1(x)&=
\begin{cases}
x\sin(\frac{1}{x}),&x\neq0,\\
0,&x=0,
\end{cases}\\
f_2(x)&=
\begin{cases}
x^2\sin(\frac{1}{x}),&x\neq0,\\
0,&x=0.
\end{cases}
\end{align*}
Quotient $\frac{f_1(t)-f_1(0)}{t}=\sin(\frac{1}{t})$ has no limit, so $f_1'(0)$ does not exist. Since $\abs{\frac{f_2(t)-f_2(0)}{t}}\leq\abs{t}$, $f_2'(0)=0$.\\
For $x\neq0$, $f_2'(x)=2x\sin(\frac{1}{x})-\cos(\frac{1}{x})$. Along $x_n=(2\pi n)^{-1}$ and $y_n=((2n+1)\pi)^{-1}$, $f_2'(x_n)=-1$ and $f_2'(y_n)=1$.\\
Thus $f_2$ is differentiable at $0$, but $f_2'$ is not continuous there.
\end{example}

\begin{definition} \label{def:local-extrema} \hlt{Local Extrema}\\
Let $(X,d)$ be a metric space, let $f:X\to\R$, and let $p\in X$.
\begin{enumerate}[label=(\roman*)]
\setlength{\itemsep}{0pt}
\item The map $f$ has a \hlt{local maximum at $p$} $\Longleftrightarrow(\exists\delta>0)(\forall q\in X)\,[d(p,q)<\delta\Longrightarrow f(q)\leq f(p)]$.
\item The map $f$ has a \hlt{local minimum at $p$} $\Longleftrightarrow(\exists\delta>0)(\forall q\in X)\,[d(p,q)<\delta\Longrightarrow f(q)\geq f(p)]$.
\end{enumerate}
\end{definition}

\begin{remark} \label{rem:local-versus-global-extrema} \hlt{Local and Global Extrema}\\
Every global extremum is local, whereas every isolated domain point is both a local maximum and minimum.\\
Derivative tests therefore require an interior interval point.\\
Continuous real maps on nonempty compact spaces attain global extrema by Theorem \ref{thm:extreme-value-theorem-compact}.
\end{remark}

\begin{theorem} \label{thm:interior-extremum-derivative-zero} \hlt{Interior Extremum Has Zero Derivative}\\
Let $a<b$, $f:(a,b)\to\R$, $x\in(a,b)$. If $f$ has a local extremum at $x$ and is differentiable there, then $f'(x)=0$.
\begin{proof}
At a local maximum, $\frac{f(t)-f(x)}{t-x}\geq0$ for $t<x$ near $x$ and $\frac{f(t)-f(x)}{t-x}\leq0$ for $t>x$ near $x$.\\
The derivative limit therefore satisfies $f'(x)\geq0$ and $f'(x)\leq0$, so $f'(x)=0$. At a local minimum, $-f$ has a local maximum.\\
The identity $(-f)'(x)=-f'(x)$ holds by Definition \ref{def:derivative-real-function}, so $f'(x)=0$ by the maximum case.\\
\hlt{Structural Commentary}: The two approach sides impose opposite sign bounds on the same derivative limit.\\
\hlt{Key Insight}: A two-sided limit compatible with both quotient signs must be zero.
\end{proof}
\end{theorem}

\begin{theorem} \label{thm:rolles-theorem-one-variable} \hlt{Rolle's Theorem}\\
Let $a<b$ and let $f:[a,b]\to\R$ be continuous on $[a,b]$ and differentiable on $(a,b)$. If $f(a)=f(b)$, then $(\exists x\in(a,b))\,f'(x)=0$.
\begin{proof}
If $f$ is constant, $f'(\frac{a+b}{2})=0$ by Definition \ref{def:derivative-real-function}. Otherwise, $[a,b]$ is compact by Theorem \ref{thm:heine-borel-rk}.\\
The map $f$ attains a maximum and a minimum by Theorem \ref{thm:extreme-value-theorem-compact}.\\
At least one extremal value differs from $f(a)=f(b)$, so its attaining point $x$ lies in $(a,b)$. Then $f'(x)=0$ by Theorem \ref{thm:interior-extremum-derivative-zero}.\\
\hlt{Structural Commentary}: Compactness supplies extrema, and equal endpoint values force a nonconstant extremum into the interior.\\
\hlt{Key Insight}: An interior attained extremum is stationary.
\end{proof}
\end{theorem}

\begin{theorem} \label{thm:cauchy-mean-value-theorem} \hlt{Cauchy Mean Value Theorem}\\
Let $a<b$, $f,g:[a,b]\to\R$ be continuous on $[a,b]$, differentiable on $(a,b)$. Then $(\exists x\in(a,b))\,\bigl(f(b)-f(a)\bigr)g'(x)=\bigl(g(b)-g(a)\bigr)f'(x)$.
\begin{proof}
Set $h(t)\defeq\bigl(f(b)-f(a)\bigr)g(t)-\bigl(g(b)-g(a)\bigr)f(t)$; then $h(b)-h(a)=0$. The map $h$ is continuous on $[a,b]$ by Theorem \ref{thm:continuous-algebra-and-components}(i).\\
It is differentiable on $(a,b)$ by Theorem \ref{thm:derivative-algebra-rules}(i),(ii). Choose $x\in(a,b)$ with $h'(x)=0$ by Theorem \ref{thm:rolles-theorem-one-variable}.\\
Since $h'(x)=\bigl(f(b)-f(a)\bigr)g'(x)-\bigl(g(b)-g(a)\bigr)f'(x)$, the claimed identity follows.\\
\hlt{Structural Commentary}: A determinant-like linear combination converts the desired identity into Rolle's endpoint condition.\\
\hlt{Key Insight}: Cancel the endpoint increments before differentiating.
\end{proof}
\end{theorem}

\begin{corollary} \label{cor:mean-value-theorem} \hlt{Mean Value Theorem}\\
Let $a<b$, $f:[a,b]\to\R$ be continuous on $[a,b]$ and differentiable on $(a,b)$. Then $(\exists x\in(a,b))\,f(b)-f(a)=(b-a)f'(x)$.
\begin{proof}
Set $g(t)\defeq t$. Then $g'(x)=1$ and $g(b)-g(a)=b-a$. Choose $x\in(a,b)$ with $f(b)-f(a)=(b-a)f'(x)$ by Theorem \ref{thm:cauchy-mean-value-theorem}.\\
\hlt{Structural Commentary}: Specializing Cauchy's auxiliary comparison to the identity map isolates the secant increment.\\
\hlt{Key Insight}: The identity map converts Cauchy's relation into slope equality.
\end{proof}
\end{corollary}

\noindent Figure \ref{fig:mean-value-secant-tangent} depicts Corollary \ref{cor:mean-value-theorem}. The tangent at an interior point $c$ is parallel to the endpoint secant.

\begin{illustrationfigure}[H]
  \centering
  \begin{tikzpicture}[
  scale=0.72,
  >=Stealth,
  every node/.style={font=\scriptsize},
  axis/.style={->, black!70, line width=0.5pt},
  curve/.style={black, line width=0.85pt},
  secant/.style={black, dashed, dash pattern=on 8pt off 6pt, line width=0.65pt},
  tangent/.style={<->, green!45!black, line width=0.75pt},
  guide/.style={black!55, dashed, line width=0.45pt},
  tick/.style={black!70, line width=0.45pt},
  point/.style={circle, fill=black, inner sep=1.45pt},
  labelbox/.style={fill=white, fill opacity=0.94, text opacity=1, inner sep=1.2pt}
]
  \draw[axis] (-0.55,0) -- (7.25,0);
  \draw[axis] (0,-0.5) -- (0,4.55);

  \coordinate (A) at (1.25,1.05);
  \coordinate (C) at (3.85,1.38);
  \coordinate (B) at (5.65,3.35);

  \draw[curve, <-]
    (0.42,1.42) .. controls (0.72,1.28) and (0.98,1.12) .. (A);
  \draw[curve]
    (A) .. controls (1.80,0.78) and (2.95,0.91) .. (C);
  \draw[curve]
    (C) .. controls (4.75,1.85) and (4.95,2.85) .. (B);
  \draw[curve, ->]
    (B) .. controls (6.05,4.10) and (6.35,4.32) .. (6.90,4.52);

  \draw[secant] (0.62,0.72) -- (6.72,3.91);
  \draw[tangent] (2.82,0.84) -- (5.10,2.03);

  \draw[guide] (C) -- (3.85,0);

  \foreach \x/\lab in {1.25/a,3.85/c,5.65/b} {
    \draw[tick] (\x,0.12) -- (\x,-0.12);
    \node[below] at (\x,-0.12) {$\lab$};
  }

  \node[point] at (A) {};
  \node[point] at (B) {};

  \node[labelbox, anchor=south east] at ($(A)+(-0.06,0.12)$) {$(a,f(a))$};
  \node[labelbox, anchor=south west] at ($(B)+(-0.02,0.08)$) {$(b,f(b))$};
  \node[labelbox, anchor=west] at (5.08,1.76)
    {\textcolor{green!45!black}{$f'(c)=\dfrac{f(b)-f(a)}{b-a}$}};
\end{tikzpicture}

  \caption{The Mean Value Theorem finds an interior tangent slope equal to the endpoint secant slope.}
  \label{fig:mean-value-secant-tangent}
\end{illustrationfigure}

\begin{corollary} \label{cor:bounded-derivative-lipschitz-one-variable} \hlt{Lipschitz Control from a Bounded Derivative}\\
Let $a<b$, let $M\geq0$, and let $f:[a,b]\to\R$ be continuous on $[a,b]$ and differentiable on $(a,b)$. If $(\forall t\in(a,b))\,\abs{f'(t)}\leq M$, then:
\begin{enumerate}[label=(\roman*)]
\setlength{\itemsep}{0pt}
\item $(\forall x,y\in[a,b])\,\abs{f(y)-f(x)}\leq M\abs{y-x}$;
\item $f$ is uniformly continuous on $[a,b]$.
\end{enumerate}
\begin{proof}
\begin{enumerate}[label=(\roman*)]
\setlength{\itemsep}{0pt}
\item If $x<y$, choose $c\in(x,y)$ with $f(y)-f(x)=(y-x)f'(c)$ by Corollary \ref{cor:mean-value-theorem}.\\
Hence $\abs{f(y)-f(x)}\leq M\abs{y-x}$. Interchanging $x,y$ covers $x>y$, and $x=y$ is immediate.
\item Given $\varepsilon>0$, set $\delta\defeq\frac{\varepsilon}{M+1}$. For $x,y\in[a,b]$, $\abs{x-y}<\delta$ implies $\abs{f(x)-f(y)}\leq M\abs{x-y}<\varepsilon$ by (i).
\end{enumerate}
\hlt{Structural Commentary}: The mean-value identity transfers the derivative bound to secants, which supplies a uniform modulus.\\
\hlt{Key Insight}: One global slope bound controls every pair of points.
\end{proof}
\end{corollary}

\begin{theorem} \label{thm:derivative-sign-monotonicity} \hlt{Derivative Sign and Monotonicity}\\
Let $a<b$ and let $f:(a,b)\to\R$ be differentiable.
\begin{enumerate}[label=(\roman*)]
\setlength{\itemsep}{0pt}
\item If $(\forall x\in(a,b))\,f'(x)\geq0$, then $f$ is monotonically increasing.
\item If $(\forall x\in(a,b))\,f'(x)=0$, then $f$ is constant.
\item If $(\forall x\in(a,b))\,f'(x)\leq0$, then $f$ is monotonically decreasing.
\end{enumerate}
\begin{proof}
Fix $x_1,x_2\in(a,b)$ with $x_1<x_2$. The map $f$ is continuous on $[x_1,x_2]$ by Theorem \ref{thm:differentiable-implies-continuous}.\\
Choose $c\in(x_1,x_2)$ with $f(x_2)-f(x_1)=(x_2-x_1)f'(c)$ by Corollary \ref{cor:mean-value-theorem}.
\begin{enumerate}[label=(\roman*)]
\setlength{\itemsep}{0pt}
\item If $f'\geq0$, then $f(x_2)-f(x_1)\geq0$.
\item If $f'=0$, then $f(x_2)-f(x_1)=0$.
\item If $f'\leq0$, then $f(x_2)-f(x_1)\leq0$.
\end{enumerate}
\hlt{Structural Commentary}: The mean-value identity transfers each derivative-sign hypothesis to every ordered pair.\\
\hlt{Key Insight}: Tangent-sign control becomes secant-sign control.
\end{proof}
\end{theorem}

\begin{theorem} \label{thm:derivatives-have-intermediate-value-property} \hlt{Derivatives Have the Intermediate Value Property}\\
Let $a<b$, $f:[a,b]\to\R$ be differentiable relative to $[a,b]$. If $\lambda\in\R$ and $\bigl(f'(a)-\lambda\bigr)\bigl(f'(b)-\lambda\bigr)<0$, then $(\exists x\in(a,b))\,f'(x)=\lambda$.
\begin{proof}
The transformation $(f,\lambda)\mapsto(-f,-\lambda)$ interchanges the endpoint orders and preserves the conclusion.\\
It therefore suffices to assume $f'(a)<\lambda<f'(b)$. Set $g(t)\defeq f(t)-\lambda t$. Then $g'(a)<0<g'(b)$ by Definition \ref{def:derivative-real-function}.\\
Choose $t_1,t_2\in(a,b)$ with $g(t_1)<g(a)$, $g(t_2)<g(b)$ by two endpoint derivative limits. Map $g$ is continuous on $[a,b]$ by Theorem \ref{thm:differentiable-implies-continuous}.\\
The interval $[a,b]$ is compact by Theorem \ref{thm:heine-borel-rk}, so $g$ attains a minimum by Theorem \ref{thm:extreme-value-theorem-compact}.\\
Neither endpoint is a minimizer, so a minimizer $x$ lies in $(a,b)$. Then $0=g'(x)=f'(x)-\lambda$ by Theorem \ref{thm:interior-extremum-derivative-zero}.\\
\hlt{Structural Commentary}: Linear tilting turns the target derivative value into an interior-minimum condition.\\
\hlt{Key Insight}: Opposite endpoint slopes force the tilted minimum away from both endpoints.
\end{proof}
\end{theorem}

\begin{perspective} \label{per:sliding-secant-proof-of-darboux} \hlt{Sliding Secants}\\
Let $f,\lambda$ be as in Theorem \ref{thm:derivatives-have-intermediate-value-property}, and fix $0<h<b-a$. Define $\sigma_h:[a,b-h]\to\R$ by $\sigma_h(x)\defeq\frac{f(x+h)-f(x)}{h}$.\\
Map $f$ is continuous by Theorem \ref{thm:differentiable-implies-continuous}. Thus $\sigma_h$ is continuous by Theorems \ref{thm:continuous-composition} and \ref{thm:continuous-algebra-and-components}(i).\\
As $h\to0^+$, $\sigma_h(a)\to f'(a)$ and $\sigma_h(b-h)\to f'(b)$. For sufficiently small $h$, $\lambda$ lies strictly between these endpoint values.\\
Choose $x\in(a,b-h)$ with $\lambda=\sigma_h(x)$ by Theorem \ref{thm:intermediate-value-theorem}. Choose $c\in(x,x+h)$ with $\lambda=f'(c)$ by Corollary \ref{cor:mean-value-theorem}.\\
\hlt{Key Insight}: Slide a fixed-width secant, then convert its slope to a tangent slope.
\end{perspective}

\noindent Figure \ref{fig:sliding-secant-darboux-derivative} illustrates the fixed-width secants in the preceding perspective.\\
Their horizontal separation remains $h$ while the base point moves.

\begin{illustrationfigure}[H]
  \centering
  \begin{tikzpicture}[
  x=0.72cm,
  y=0.67cm,
  >=Stealth,
  every node/.style={font=\scriptsize},
  axis/.style={->, black!55, line width=0.45pt},
  curve/.style={black!85, line width=0.9pt},
  secant/.style={black!65, line width=0.8pt},
  guide/.style={black!22, dashed, line width=0.38pt},
  tick/.style={black!55, line width=0.45pt},
  point/.style={circle, fill=black, inner sep=1.15pt},
  hbrace/.style={<->, black!55, line width=0.45pt}
]
  \draw[axis] (-0.35,0) -- (8.45,0) node[right] {$t$};
  \draw[axis] (0,-0.25) -- (0,4.15);

  \coordinate (A) at (0.75,3.40);
  \coordinate (B) at (1.65,2.00);
  \coordinate (C) at (3.75,1.42);
  \coordinate (D) at (4.65,1.56);
  \coordinate (E) at (6.05,2.65);
  \coordinate (F) at (6.95,3.58);

  \draw[curve]
    (0.30,3.90) .. controls (0.55,3.65) and (0.62,3.52) .. (A)
    .. controls (1.00,2.72) and (1.28,2.18) .. (B)
    .. controls (2.25,1.08) and (3.12,1.05) .. (C)
    .. controls (4.05,1.35) and (4.32,1.48) .. (D)
    .. controls (5.15,1.72) and (5.55,2.18) .. (E)
    .. controls (6.35,3.00) and (6.55,3.28) .. (F)
    .. controls (7.35,3.82) and (7.70,3.72) .. (8.05,3.52);

  \foreach \p/\q in {A/B, C/D, E/F} {
    \draw[secant] (\p) -- (\q);
    \node[point] at (\p) {};
    \node[point] at (\q) {};
  }

  \foreach \x in {0.75,1.65,3.75,4.65,6.05,6.95} {
    \draw[guide] (\x,0) -- (\x,3.75);
    \draw[tick] (\x,0.09) -- (\x,-0.09);
  }

  \foreach \x/\lab in {0.75/{x_1},3.75/{x}} {
    \node[below] at (\x,-0.08) {$\lab$};
  }
  \node[below] at (6.95,-0.08) {$x_2$};

  \foreach \xl/\xr in {0.75/1.65,3.75/4.65,6.05/6.95} {
    \draw[hbrace] (\xl,-0.58) -- node[below] {$h$} (\xr,-0.58);
  }

  \draw[->, black!45] (2.15,0.52) -- (3.18,0.52);
  \draw[->, black!45] (5.05,0.52) -- (5.62,0.52);
\end{tikzpicture}

  \caption{Fixed-width secants slide across the graph while retaining horizontal separation $h$.}
  \label{fig:sliding-secant-darboux-derivative}
\end{illustrationfigure}

\begin{corollary} \label{cor:derivatives-no-simple-discontinuities} \hlt{Derivatives Have No Simple Discontinuities}\\
Let $a<b$, $f:[a,b]\to\R$ be differentiable relative to $[a,b]$. Then $f'$ has no simple discontinuity in $(a,b)$.
\begin{proof}
Fix $c\in(a,b)$, let $L$ be either existing one-sided limit $f'(c-)$ or $f'(c+)$. If $L\neq f'(c)$, choose $\lambda$ strictly between them.\\
The function $f'$ remains strictly on the $L$-side of $\lambda$ in some one-sided neighborhood of $c$ by Definition \ref{def:one-sided-limits-and-discontinuities}.\\
Fix $x$ in that neighborhood. The endpoint derivatives of restriction to the segment with endpoints $x,c$ are $f'(x)$ and $f'(c)$.\\
Choose $z$ strictly between $x$ and $c$ with $f'(z)=\lambda$ by Theorem \ref{thm:derivatives-have-intermediate-value-property}, a contradiction.\\
Thus every existing one-sided limit of $f'$ at $c$ equals $f'(c)$. If $c$ a simple discontinuity, both one-sided limits exist by Definition \ref{def:one-sided-limits-and-discontinuities}.\\
Then $f'(t)\to f'(c)$ by Proposition \ref{prop:two-sided-limit-from-one-sided-limits}. Thus $f'$ would be continuous at $c$ by Theorem \ref{thm:continuity-as-limit-at-limit-point}.\\
\hlt{Structural Commentary}: Darboux interpolation forces each existing one-sided derivative limit to equal the central derivative.\\
\hlt{Key Insight}: Any gap from $f'(c)$ would force the forbidden value arbitrarily nearby.
\end{proof}
\end{corollary}

\begin{theorem} \label{thm:inverse-function-derivative-one-variable} \hlt{Derivative of an Inverse Function}\\
Let $I\subseteq\R$ be an open interval, $f:I\to\R$ be injective, $x_0\in I$. Set $y_0\defeq f(x_0)$ and $g\defeq f^{-1}:f(I)\to I$.\\
If $f'(x_0)$ exists and is nonzero and $g$ is continuous at $y_0$, then $g'(y_0)=\frac{1}{f'(x_0)}$ relative to $f(I)$.
\begin{proof}
Choose $(x_n)\subseteq I\setminus\{x_0\}$ with $x_n\to x_0$. Then $f(x_n)\to y_0$ by Theorem \ref{thm:differentiable-implies-continuous}, and $f(x_n)\neq y_0$ by injectivity.\\
Thus $y_0\in f(I)^{\operatorname{acc}}$. For $y\in f(I)\setminus\{y_0\}$, set $x\defeq g(y)$. Then $\frac{g(y)-g(y_0)}{y-y_0}=\frac{x-x_0}{f(x)-f(x_0)}=\left(\frac{f(x)-f(x_0)}{x-x_0}\right)^{-1}$.\\
As $y\to y_0$, continuity of $g$ implies $x=g(y)\to x_0$. The reciprocal therefore tends to $\frac{1}{f'(x_0)}$ by Theorem \ref{thm:function-limit-algebra}(iii).\\
\hlt{Structural Commentary}: Inverse continuity transports the input limit, while injectivity makes the matched quotients reciprocal.\\
\hlt{Key Insight}: Exchange input and output increments before taking the limit.
\end{proof}
\end{theorem}

\begin{corollary} \label{cor:one-dimensional-inverse-function-theorem} \hlt{One-Dimensional Inverse Function Theorem}\\
Let $a<b$, $c<d$, $f:(a,b)\to(c,d)$ be a differentiable bijection. Assume $(\forall x\in(a,b))\,f'(x)\neq0$.
\begin{enumerate}[label=(\roman*)]
\setlength{\itemsep}{0pt}
\item The map $f$ is a homeomorphism.
\item The inverse $g\defeq f^{-1}$ is differentiable and satisfies $(\forall y\in(c,d))\,g'(y)=\frac{1}{f'(g(y))}$.
\end{enumerate}
\begin{proof}
The derivative $f'$ has constant sign by Theorem \ref{thm:derivatives-have-intermediate-value-property} and the nonvanishing hypothesis. Fix $x_1,x_2\in(a,b)$ with $x_1<x_2$.\\
Choose $\xi\in(x_1,x_2)$ with $f(x_2)-f(x_1)=(x_2-x_1)f'(\xi)$ by Corollary \ref{cor:mean-value-theorem}. Hence $f$ is strictly monotone.
\begin{enumerate}[label=(\roman*)]
\setlength{\itemsep}{0pt}
\item The inverse $g:(c,d)\to(a,b)$ is monotone, and its image is the interval $(a,b)$.\\
Thus $g$ is continuous by Theorem \ref{thm:monotone-continuity-interval-image}. The map $f$ is continuous by Theorem \ref{thm:differentiable-implies-continuous}, so $f$ is a homeomorphism.
\item Fix $y\in(c,d)$ and set $x\defeq g(y)$. Then $g'(y)=\frac{1}{f'(x)}=\frac{1}{f'(g(y))}$ by Theorem \ref{thm:inverse-function-derivative-one-variable}.
\end{enumerate}
\hlt{Structural Commentary}: Darboux fixes the derivative sign, monotonicity controls the inverse, and the local quotient formula finishes.\\
\hlt{Key Insight}: A nonvanishing derivative prevents both turning and reciprocal-quotient failure.
\end{proof}
\end{corollary}

\begin{theorem} \label{thm:lhospital-one-sided} \hlt{L'Hospital's Rule}\\
Let $-\infty\leq a<b\leq+\infty$, $A\in\overline{\R}$, $f,g:(a,b)\to\R$ be differentiable. Assume $(\forall x\in(a,b))\,g'(x)\neq0$.
\begin{enumerate}[label=(\roman*)]
\setlength{\itemsep}{0pt}
\item If $\frac{f'(x)}{g'(x)}\to A$ and $f(x),g(x)\to0$ as $x\to a^+$, then $\frac{f(x)}{g(x)}\to A$.
\item If $\frac{f'(x)}{g'(x)}\to A$ and $\abs{g(x)}\to+\infty$ as $x\to a^+$, then $\frac{f(x)}{g(x)}\to A$.
\item If $\frac{f'(x)}{g'(x)}\to A$ as $x\to b^-$, the conclusions in (i)--(ii) hold with $b^-$ in place of $a^+$.
\end{enumerate}
\begin{proof}
The derivative $g'$ has constant sign by Theorem \ref{thm:derivatives-have-intermediate-value-property} and the nonvanishing hypothesis.\\
For $r<s$, choose $t\in(r,s)$ with $g(s)-g(r)=(s-r)g'(t)$ by Corollary \ref{cor:mean-value-theorem}. Hence $g$ is strictly monotone.\\
For $a<r<s<b$, set $Q(r,s)\defeq\frac{f(s)-f(r)}{g(s)-g(r)}$. Choose $t\in(r,s)$ with $Q(r,s)=\frac{f'(t)}{g'(t)}$ by Theorem \ref{thm:cauchy-mean-value-theorem}.\\
For (i)--(ii), every eventual order bound on $\frac{f'}{g'}$ near $a$ also holds for $Q(r,s)$ when $r,s$ are sufficiently near $a$.
\begin{enumerate}[label=(\roman*)]
\setlength{\itemsep}{0pt}
\item After restricting to a smaller right-neighborhood of $a$, $g\neq0$ by strict monotonicity and $g(x)\to0$.\\
Fix $s$ there and let $r\to a^+$. Then $Q(r,s)\to\frac{f(s)}{g(s)}$.\\
Each tightened eventual bound for $\frac{f'}{g'}$ therefore transfers to $\frac{f(s)}{g(s)}$. Thus $\frac{f(x)}{g(x)}\to A$ by Definition \ref{def:extended-real-function-limits}.
\item The limit $\abs{g(r)}\to+\infty$ makes $g(r)\neq0$ near $a$. Fix $s$ sufficient near $a$.\\
For $r\in(a,s)$, $\frac{f(r)}{g(r)}=Q(r,s)\frac{g(r)-g(s)}{g(r)}+\frac{f(s)}{g(r)}$.\\
As $r\to a^+$, the second factor tends to $1$ and the last term tends to $0$.\\
For finite $A$, use a half-width interval; for $A=\pm\infty$, use a farther ray threshold.\\
Each tightened bound remains inside prescribed neighborhood after these perturbations. Thus $\frac{f(r)}{g(r)}\to A$ by Definition \ref{def:extended-real-function-limits}.
\item Define $F(u)\defeq f(-u)$ and $G(u)\defeq g(-u)$ on $(-b,-a)$. Then $\frac{F'(u)}{G'(u)}=\frac{f'(-u)}{g'(-u)}$.\\
The left-endpoint conclusions for $F,G$ hold by (i)--(ii) and translate to the stated limits as $x\to b^-$.
\end{enumerate}
\hlt{Structural Commentary}: Cauchy's quotient identity transfers derivative bounds; endpoint limits remove or dominate the base terms.\\
\hlt{Key Insight}: Control every two-point quotient before sending one endpoint to the boundary.
\end{proof}
\end{theorem}

\begin{definition} \label{def:higher-derivatives} \hlt{Higher Derivatives}\\
Let $I\subseteq\R$ be an interval and let $f:I\to\R$. Set $f^{(0)}\defeq f$.\\
For $n\in\N$, define $f^{(n)}\defeq\bigl(f^{(n-1)}\bigr)'$ relative to $I$ whenever this derivative exists on $I$.\\
The map $f$ is \hlt{$n$ times differentiable} on $I$ $\Longleftrightarrow(\forall j\in\{1,\ldots,n\})\,f^{(j)}$ exists on $I$.
\end{definition}

\begin{definition} \label{def:smoothness-classes-one-variable} \hlt{Smoothness Classes}\\
Let $I\subseteq\R$ be an interval.
\begin{enumerate}[label=(\roman*)]
\setlength{\itemsep}{0pt}
\item $C^0(I)$ is the class of continuous maps $f:I\to\R$.
\item For $r\geq1$, $C^r(I)$ is the class of $r$ times differentiable maps $f:I\to\R$ with $f^{(r)}$ continuous.
\item $C^\infty(I)\defeq\bigcap\limits_{r=0}^{\infty}C^r(I)$ is the class of \hlt{smooth} maps on $I$.
\end{enumerate}
\end{definition}

\begin{theorem} \label{thm:higher-order-leibniz-product-rule} \hlt{Higher-Order Leibniz Product Rule}\\
Let $a<b$, $r\in\N$, $f,g:(a,b)\to\R$ be $r$ times differentiable. Then $(\forall x\in(a,b))\,(fg)^{(r)}(x)=\sum\limits_{k=0}^{r}\binom{r}{k}f^{(k)}(x)g^{(r-k)}(x)$.
\begin{proof}
For $r=1$, formula holds by Theorem \ref{thm:derivative-algebra-rules}(ii). Assume it holds at order $r$, $f,g$ be $(r+1)$ times differentiable.\\
The differentiated identity is $(fg)^{(r+1)}=\sum\limits_{k=0}^{r}\binom{r}{k}f^{(k+1)}g^{(r-k)}+\sum\limits_{k=0}^{r}\binom{r}{k}f^{(k)}g^{(r-k+1)}$ by Theorem \ref{thm:derivative-algebra-rules}(ii).\\
After reindexing the first sum by $j=k+1$, the coefficient of $f^{(j)}g^{(r+1-j)}$ is $\binom{r}{j-1}+\binom{r}{j}=\binom{r+1}{j}$.\\
The endpoint coefficients at $j=0,r+1$ are $1$, so the order-$(r+1)$ formula follows.\\
\hlt{Structural Commentary}: Differentiation creates two shifted sums, and reindexing aligns their common derivative orders.\\
\hlt{Key Insight}: Pascal's identity merges the two coefficient contributions.
\end{proof}
\end{theorem}

\begin{theorem} \label{thm:taylor-theorem-lagrange-remainder} \hlt{Taylor's Theorem with Lagrange Remainder}\\
Let $a<b$, $n\in\N$, $f:[a,b]\to\R$, $f^{(n-1)}$ continuous on $[a,b]$, $\exists$ $f^{(n)}$ on $(a,b)$. For distinct $\alpha,\beta\in[a,b]$, $J_{\alpha,\beta}\defeq(\min\{\alpha,\beta\},\max\{\alpha,\beta\})$.\\
Define $P(t)\defeq\sum\limits_{k=0}^{n-1}\frac{f^{(k)}(\alpha)}{k!}(t-\alpha)^k$. Then $(\exists x\in J_{\alpha,\beta})\,f(\beta)=P(\beta)+\frac{f^{(n)}(x)}{n!}(\beta-\alpha)^n$.
\begin{proof}
\begin{enumerate}[label=(\roman*)]
\setlength{\itemsep}{0pt}
\item \hlt{Normalization}: Set $M\defeq\frac{f(\beta)-P(\beta)}{(\beta-\alpha)^n}$ and $G(t)\defeq f(t)-P(t)-M(t-\alpha)^n$.
\item \hlt{Endpoint cancellation}: For $0\leq j\leq n-1$, $P^{(j)}(\alpha)=f^{(j)}(\alpha)$, so $G^{(j)}(\alpha)=0$. Identity $G(\beta)=0$ holds by definition of $M$.
\item \hlt{Iterated Rolle step}: Set $c_0\defeq\beta$ and $J(u,v)\defeq(\min\{u,v\},\max\{u,v\})$.\\
For $0\leq j<n$, $G^{(j)}$ is continuous on $[a,b]$ by the regularity hypotheses and Theorem \ref{thm:differentiable-implies-continuous}.\\
It is differentiable on $(a,b)$ by the regularity hypotheses. Fix $0\leq j<n$ with $G^{(j)}(c_j)=G^{(j)}(\alpha)=0$.\\
Choose $c_{j+1}\in J(\alpha,c_j)$ with $G^{(j+1)}(c_{j+1})=0$ by Theorem \ref{thm:rolles-theorem-one-variable}. This recursion reaches $c_n\in J_{\alpha,\beta}$.
\item \hlt{Remainder identification}: For $x\defeq c_n$, $0=G^{(n)}(x)=f^{(n)}(x)-n!M$. Thus $M=\frac{f^{(n)}(x)}{n!}$, which is the claimed formula.
\end{enumerate}
\hlt{Structural Commentary}: Endpoint cancellation feeds an $n$-step Rolle chain whose last zero identifies the normalized remainder.\\
\hlt{Key Insight}: Force the error to vanish once at $\beta$ and $n$ times at $\alpha$.
\end{proof}
\end{theorem}

\begin{definition} \label{def:vector-valued-derivative-one-variable} \hlt{Vector-Valued Derivative}\\
Let $k\in\N$, $X\subseteq\R$, $x_0\in X^{\operatorname{acc}}$, $\mathbf{f}:X\to\R^k$. For $\mathbf{L}\in\R^k$ and $t\in X\setminus\{x_0\}$, set $\mathbf{q}_{\mathbf{L}}(t)\defeq\frac{\mathbf{f}(t)-\mathbf{f}(x_0)}{t-x_0}-\mathbf{L}$.\\
Then $\mathbf{f}'(x_0)=\mathbf{L}$ relative to $X$ means $\lim\limits_{t\to x_0}\norm{\mathbf{q}_{\mathbf{L}}(t)}_2=0$.
\end{definition}

\begin{proposition} \label{prop:vector-derivative-componentwise} \hlt{Componentwise Characterization of the Vector Derivative}\\
Let $k\in\N$, $X\subseteq\R$, $\mathbf{f}=(f_1,\ldots,f_k):X\to\R^k$, $x_0\in X^{\operatorname{acc}}$. Let $\mathbf{L}=(L_1,\ldots,L_k)\in\R^k$; all derivatives are relative to $X$.\\
Then $\mathbf{f}'(x_0)=\mathbf{L}$ $\Longleftrightarrow(\forall i\in\{1,\ldots,k\})\,f_i'(x_0)=L_i$.
\begin{proof}
For $t\in X\setminus\{x_0\}$, set $\mathbf{q}(t)\defeq\frac{\mathbf{f}(t)-\mathbf{f}(x_0)}{t-x_0}-\mathbf{L}$.
\begin{enumerate}[label=\roman*.]
\setlength{\itemsep}{0pt}
\item[$(\Rightarrow)$] If $\norm{\mathbf{q}(t)}_2\to0$, then $(\forall i\in\{1,\ldots,k\})\,\abs{q_i(t)}\leq\norm{\mathbf{q}(t)}_2\to0$. Hence $(\forall i\in\{1,\ldots,k\})\,f_i'(x_0)=L_i$ by Definition \ref{def:derivative-real-function}.
\item[$(\Leftarrow)$] If $(\forall i\in\{1,\ldots,k\})\,q_i(t)\to0$, then $\norm{\mathbf{q}(t)}_2\leq\sqrt{k}\max\limits_{1\leq i\leq k}\abs{q_i(t)}\to0$. Hence $\mathbf{f}'(x_0)=\mathbf{L}$ by Definition \ref{def:vector-valued-derivative-one-variable}.
\end{enumerate}
\hlt{Structural Commentary}: Coordinate bounds compare the vector quotient error with its finitely many scalar errors.\\
\hlt{Key Insight}: Finite-dimensional norm convergence is exactly coordinatewise convergence.
\end{proof}
\end{proposition}

\begin{theorem} \label{thm:vector-valued-mean-value-inequality} \hlt{Vector-Valued Mean Value Inequality}\\
Let $k\in\N$, $a<b$, $\mathbf{f}:[a,b]\to\R^k$ be continuous on $[a,b]$, differentiable on $(a,b)$. Then $(\exists x\in(a,b))\,\norm{\mathbf{f}(b)-\mathbf{f}(a)}_2\leq(b-a)\norm{\mathbf{f}'(x)}_2$.
\begin{proof}
Set $\mathbf{z}\defeq\mathbf{f}(b)-\mathbf{f}(a)$ and $\varphi(t)\defeq\mathbf{z}\cdot\mathbf{f}(t)$. Map $\varphi$ is continuous on $[a,b]$ by Theorem \ref{thm:continuous-algebra-and-components}(ii),(iii).\\
For $t\in(a,b)$ and $s\in[a,b]\setminus\{t\}$, $\frac{\varphi(s)-\varphi(t)}{s-t}=\mathbf{z}\cdot\frac{\mathbf{f}(s)-\mathbf{f}(t)}{s-t}$. Thus $\varphi'(t)=\mathbf{z}\cdot\mathbf{f}'(t)$ by Theorem \ref{thm:function-limit-algebra}(iv).\\
Choose $x\in(a,b)$ with $\varphi(b)-\varphi(a)=(b-a)\mathbf{z}\cdot\mathbf{f}'(x)$ by Corollary \ref{cor:mean-value-theorem}. The identity $\varphi(b)-\varphi(a)=\norm{\mathbf{z}}_2^2$ holds.\\
Hence $\norm{\mathbf{z}}_2^2\leq(b-a)\norm{\mathbf{z}}_2\norm{\mathbf{f}'(x)}_2$ by Theorem \ref{thm:cauchy-schwarz-inequality}(i). If $\mathbf{z}\neq\mathbf{0}$, cancellation proves the claim; if $\mathbf{z}=\mathbf{0}$, every $x\in(a,b)$ satisfies it.\\
\hlt{Structural Commentary}: Scalar projection creates a mean-value identity, and Cauchy--Schwarz converts it back to norms.\\
\hlt{Key Insight}: Project along the endpoint displacement itself.
\end{proof}
\end{theorem}

\begin{proofmethod} \label{meth:one-variable-differentiation-proof-patterns} \hlt{Differentiation Proof Patterns}
\begin{enumerate}[label=(\roman*)]
\setlength{\itemsep}{0pt}
\item \hlt{Domain}: Bind $f:X\to\R$ and $x_0\in X^{\operatorname{acc}}$; isolated points admit no derivative by Definition \ref{def:derivative-real-function}.
\item \hlt{Definition}: Factor $f(t)-f(x)=(t-x)Q_x(t)$ with $Q_x(t)\to f'(x)$.
\item \hlt{Approximation}: Replace the quotient limit by a first-order error bound by Proposition \ref{prop:one-variable-first-order-approximation}.
\item \hlt{Algebra}: Establish exact difference-quotient identities before passing to limits through continuity.
\item \hlt{Mean values}: Construct $h$ with $h(a)=h(b)$; an interior stationary point then exists by Theorem \ref{thm:rolles-theorem-one-variable}.
\item \hlt{Taylor}: Set $G=f-P-M(t-\alpha)^n$ and iterate Rolle's theorem through successive derivatives.
\end{enumerate}
\hlt{Key Insight}: Normalize local errors or engineer endpoint cancellation before invoking earlier results.
\end{proofmethod}
